% =====================================================================
%  THÈME 2 — Conversions en quantité de matière
%  Niveau FACILE — Exercices
% =====================================================================
\documentclass[11pt,a4paper]{article}
\usepackage[utf8]{inputenc}
\usepackage[T1]{fontenc}
\usepackage{lmodern}
\usepackage[french]{babel}
\usepackage{amsmath,amssymb,mathtools}
\usepackage{icomma}
\usepackage[version=4]{mhchem}
\usepackage{siunitx}
\sisetup{output-decimal-marker={,},per-mode=symbol}
\DeclareSIUnit\Molar{\mole\per\litre}
\DeclareSIUnit\atm{atm}
\usepackage{xcolor}
\usepackage{tikz}
\usepackage[most]{tcolorbox}
\usepackage{enumitem}
\usepackage{array}
\usepackage{fancyhdr}
\usepackage[a4paper,margin=2.1cm,headheight=26pt,headsep=8pt]{geometry}

\definecolor{primary}{RGB}{142,93,168}
\definecolor{primarydark}{RGB}{82,45,108}
\definecolor{excol}{RGB}{0,135,90}
\newcommand{\gmol}{\si{\gram\per\mole}}
\newcommand{\Vm}{V_{\mathrm m}}

\pagestyle{fancy}\fancyhf{}
\fancyhead[L]{\small\textcolor{primarydark}{\textbf{Fiche 5} — Thème 2 : Conversions}}
\fancyhead[R]{\small\textcolor{excol}{\textsc{Facile}}}
\fancyfoot[C]{\small\thepage}
\renewcommand{\headrulewidth}{0.8pt}
\renewcommand{\headrule}{\hbox to\headwidth{\color{primary}\leaders\hrule height \headrulewidth\hfill}}

\newtcolorbox[auto counter]{exo}[1][]{enhanced,breakable,boxrule=1pt,arc=3pt,
  left=3mm,right=3mm,top=2.4mm,bottom=2.4mm,colback=excol!5,colframe=excol,
  fonttitle=\bfseries,coltitle=white,
  title={Exercice~\thetcbcounter\ \ifx\relax#1\relax\relax\else\textmd{--- \itshape #1}\fi}}

\setlength{\parindent}{0pt}
\setlength{\parskip}{3pt}

\begin{document}

\begin{center}
{\LARGE\bfseries\color{primarydark}Thème 2 — Conversions}\\[1pt]
{\large\color{excol}\textbf{Niveau Facile}}\\[2pt]
{\itshape Objectif : utiliser \emph{une seule} porte à la fois vers la mole.}\\
{\color{primary}\rule{0.6\linewidth}{1pt}}
\end{center}

\textbf{Données (masses molaires atomiques, en \gmol) :}
\ce{H}=1,0 ; \ce{C}=12,0 ; \ce{O}=16,0 ; \ce{Na}=23,0 ; \ce{S}=32,0 ; \ce{Ca}=40,0.\quad
$N_A=\SI{6,02e23}{}$ ; \; $\Vm=\SI{22,4}{\litre\per\mole}$ (TPN).

\begin{exo}[calculer une masse molaire]
Calcule la masse molaire de :\quad
\textbf{a)} l'acide sulfurique \ce{H2SO4} ;\quad
\textbf{b)} le carbonate de calcium \ce{CaCO3}.\\
\emph{(Additionne les masses molaires atomiques de tous les atomes de la formule.)}
\end{exo}

\begin{exo}[masse $\leftrightarrow$ mole]
\textbf{a)} Combien de moles y a-t-il dans \SI{20,0}{\gram} d'hydroxyde de sodium \ce{NaOH}
($M=\SI{40,0}{\gram\per\mole}$) ?\\
\textbf{b)} Quelle masse de glucose \ce{C6H12O6} ($M=\SI{180}{\gram\per\mole}$) faut-il peser pour
disposer de \SI{0,25}{\mole} ?
\end{exo}

\begin{exo}[volume d'un gaz aux TPN]
Aux conditions TPN ($\Vm=\SI{22,4}{\litre\per\mole}$) :\\
\textbf{a)} quel volume occupent \SI{0,25}{\mole} de dioxygène gazeux ?\\
\textbf{b)} combien de moles de gaz contient un volume de \SI{44,8}{\litre} ?
\end{exo}

\begin{exo}[quantité de matière en solution]
On prélève \SI{250}{\milli\litre} d'une solution de chlorure de sodium de concentration
$C=\SI{0,20}{\Molar}$.\\
\textbf{a)} Convertis d'abord le volume en litres.\quad
\textbf{b)} Combien de moles de \ce{NaCl} contient ce prélèvement ?
\end{exo}

\begin{exo}[compter les entités avec $N_A$]
\textbf{a)} Combien de molécules y a-t-il dans \SI{0,10}{\mole} d'eau ?\\
\textbf{b)} À combien de moles correspond un échantillon de $\SI{3,01e23}{}$ molécules ?
\end{exo}

\end{document}
